\chapter{Movements of Earth}
\section{Get Rolling}
\textbf{Do you think it will be morning in the whole world when it is morning in India? Provide reasons for your answer and share your thoughts in the class.}
\section{\textcolor{red}{What to Know}}
Earth's motion causes the formation of day and night as well as change of seasons.
\section{I Explore}
Earth is a spherical in shape. It is slightly bulged at the Equator and flattened at the poles. The Equator, an \textcolor{red}{imaginary} line, divides Earth into two equal parts.
\par
\vspace{0.1in}
The northern half of Earth is called the Northern Hemisphere and the southern half is called the Southern Hemisphere. The word 'hemisphere' was derived from the Greek word 'hemisphairion'. It means one half of a spherical body.
\par
Approximately 67.3 percent of Earth's landmasses and around 60.7 percent of its water bodies are situated in the Northern Hemisphere. It is estimated that around 87 percent of the world's population resides in the Northern Hemisphere. In contrast, the Southern Hemisphere has only about 32.7 percent of Earth's landmasses is located in the Southern Hemisphere.
\section{\textcolor{red}{Movements of Earth}}
In the ancient times, people believed that Earth is \textcolor{red}{stationary} and heavenly bodies move around it.
\par
\vspace{0.1in}
About 2300 years ago, Aristarchus of Samos, a Greek astronomer, found out that Earth rotates on its axis and goes around the Sun.
\par
\vspace{0.1in}
Later, in the 1500s Nicolaus Copernicus, a Polish astronomer, also found out about Earth's daily movement on its axis and its yearly movement around the Sun. Thus, there are two movements of Earth that happens at the same time.
\begin{enumerate}
    \item Movement on its axis
    \vspace{0.1in}
    \item Movement around the Sun.
\end{enumerate}
\par
\vspace{0.1in}
To understand these movements, let us begin by understanding what an axis is.